%************************************************
\chapter{State of the art}\label{cap:sota}
%************************************************

% switch between these depending on page size, or modify directly
% same for the rest of the chapters
\tikz[remember picture,overlay] \node[opacity=0.3,inner sep=0pt] at (current page.center){\includegraphics[width=\paperwidth,height=\paperheight]{./Figures/cover/sesshu_1.jpg}};
%\tikz[remember picture,overlay] \node[opacity=0.3,inner sep=0pt] at ([yshift=3cm,xshift=2cm]current page.center){\includegraphics[width=\paperwidth,height=\paperheight]{./Figures/cover/sesshu_1.jpg}};
\clearpage 

Lorem ipsum

% INVESTIGA SI HAY TRABAJOS SIMILARES EN ESPAÑITAAA
% keremos primera letra en mayúsucula tocha etc?

\section{Traffic signal control}


\citep{Robertson1991}
Where adjacent junctions in a network of urban road are
less than a mile apart and are controlled by traffic signals,
major benefits can be obtained by installing an urban traffic
control (UTC) system to coordinate the operation of the signals.
Such systems use a central computer to control the signal
sequences and to monitor their operation; the signals and computer are usually connected by voice grade data transmission
lines. The method of controlling the signal sequences can have a
significant effect on levels of urban congestion, on fuel economy, and on exhaust emissions. This is the main topic of this
paper. 
The first UTC systems came into operation in the mid-l960’s,
and since then their use has been growing rapidly.

\section{Modelling the behaviour of pedestrian and cyclists}

\section{Reinforcement learning in traffic signal control}

Lorem ipsum